\documentclass{notaaula}

        \titulonota{ENEM: ondas, radiação e tecnologia}
        \creditos{Turma 3A · Ciências da Natureza para o ENEM · Outubro/2026 · Prof. Josué Cavalcante}

        \begin{document}
        \cabecalho

        \begin{sobre}
        Bloco Tecnologia, Sociedade e Intervenção: ondas, espectro eletromagnético, comunicação, saúde e leitura de evidências.
        \end{sobre}

        \section{Relações que organizam o tema}

\begin{copiar}{Duas equações-chave}
\[
  v=\lambda f,
  \qquad E_{\mathrm{fóton}}=hf.
\]
A primeira vale para qualquer onda periódica. A segunda relaciona frequência e energia de um fóton.
No mesmo meio, mudar $f$ altera $\lambda$; para ondas eletromagnéticas no vácuo, $v=c$.
\atencao{Energia do fóton depende da frequência; intensidade depende da quantidade de energia que atravessa uma área por tempo.}
\end{copiar}

\section{Espectro eletromagnético}

\begin{copiar}{Frequência crescente}
\[
  \begin{gathered}
  \text{rádio}\rightarrow\text{micro-ondas}\rightarrow\text{infravermelho}\\
  \rightarrow\text{visível}\rightarrow\text{UV}\rightarrow\text{X}\rightarrow\text{gama}.
  \end{gathered}
\]
À direita, frequência e energia por fóton aumentam; o comprimento de onda diminui. Radiofrequência,
micro-ondas e luz visível são não ionizantes. UV mais energético, raios X e gama podem ionizar matéria.
Risco biológico depende de energia, dose, duração e tecido exposto.
\end{copiar}

\begin{exemplo}[Item comentado]
Uma emissora opera em $100\un{MHz}=1\times10^8\un{Hz}$. No ar,
\[
  \lambda=\frac{c}{f}=\frac{3\times10^8}{1\times10^8}
  =\resultado{3\un{m}}.
\]
A conversão de prefixo é parte central do item: mega significa $10^6$.
\end{exemplo}

\section{Tecnologias e fenômenos}

\begin{copiar}{Associe mecanismo e uso}
\begin{itens}
\item Difração: ondas contornam obstáculos; efeito cresce quando dimensão e $\lambda$ são comparáveis.
\item Reflexão: radar, ultrassom e formação de eco.
\item Refração: lentes, fibras e desvio ao mudar de meio.
\item Polarização: telas, filtros e redução de reflexos.
\item Interferência: franjas, filmes finos e cancelamento de ruído.
\end{itens}
No ultrassom médico a onda é mecânica, não radiação eletromagnética ionizante.
\end{copiar}

\section{Leitura crítica de risco}

\begin{copiar}{Dose e evidência}
Uma boa conclusão distingue perigo de risco. Perigo é a capacidade de causar dano; risco combina
perigo com exposição. Compare grupos, dose, tempo, controle experimental e tamanho da amostra antes
de atribuir causalidade. Correlação isolada não prova mecanismo causal.
\end{copiar}

\begin{exemplo}[Comunicação óptica]
Fibras ópticas guiam luz por reflexão interna total. A portadora tem frequência alta e permite grande
taxa de informação; repetidores e conversores compensam atenuação. O princípio óptico não significa
que energia seja transmitida sem perdas.
\end{exemplo}

\begin{dica}
Em itens com espectro, marque primeiro a direção de crescimento de $f$. Depois use
$c=\lambda f$ e $E=hf$ para evitar inverter duas relações ao mesmo tempo.
\end{dica}

\begin{exercicios}[Mini-simulado ENEM]
\begin{questoes}
\item No vácuo, dobrar a frequência de uma onda eletromagnética faz $\lambda$
  \begin{alternativas}\item dobrar\item cair à metade\item permanecer\item quadruplicar\end{alternativas}
\item A radiação de maior energia por fóton é
  \begin{alternativas}\item rádio\item infravermelho\item visível\item raios X\end{alternativas}
\item O ultrassom médico é
  \begin{alternativas}\item onda mecânica\item onda eletromagnética ionizante\item corrente contínua\item partícula alfa\end{alternativas}
\item Óculos polarizados reduzem reflexos porque
  \begin{alternativas}\item aumentam todas as frequências\item selecionam uma direção de oscilação\item produzem som\item ionizam o ar\end{alternativas}
\item Uma onda de $600\un{MHz}$ tem, no vácuo, comprimento de onda
  \begin{alternativas}\item $\dec{0,05}\un{m}$\item $\dec{0,50}\un{m}$\item $5\un{m}$\item $50\un{m}$\end{alternativas}
\item A difração é mais intensa quando
  \begin{alternativas}\item a abertura é comparável a $\lambda$\item $\lambda=0$\item não há onda\item a frequência é sempre audível\end{alternativas}
\item Explique por que ``não ionizante'' não significa automaticamente ``sem qualquer risco''.
\item Um estudo observa correlação entre exposição e efeito. Que evidência adicional ajudaria a sustentar causalidade?
\end{questoes}
\gabarito{1-b; 2-d; 3-a; 4-b; 5-b; 6-a.}
\end{exercicios}


\end{document}
